% Generated draft: requires mathematical and source review.
\documentclass[11pt]{article}
\usepackage[margin=25mm]{geometry}
\usepackage{fontspec}
\setmainfont{Latin Modern Roman}
\usepackage{amsmath,amssymb,mathtools,unicode-math}
\setmathfont{Latin Modern Math}
\usepackage{enumitem,needspace,xurl,hyperref}
\setlength{\parindent}{0pt}
\setlength{\parskip}{5pt}
\setlength{\emergencystretch}{3em}
\newcommand{\N}{\mathbb{N}}
\newcommand{\Z}{\mathbb{Z}}
\newcommand{\Q}{\mathbb{Q}}
\newcommand{\R}{\mathbb{R}}
\newcommand{\C}{\mathbb{C}}
\newcommand{\F}{\mathbb{F}}
\newcommand{\A}{\mathbb{A}}
\newcommand{\PP}{\mathbb{P}}
\newcommand{\E}{\mathbb{E}}
\newcommand{\eps}{\varepsilon}
\newcommand{\dd}{\,\mathrm{d}}
\begin{document}
\section*{UnsolvedMath ID: 30005082}
\textbf{Source problem code:} catalogue entry
\subsection{Definitions and mathematical statement}
For integers \(k,n\) with \(2<k\leq n\), let \(\operatorname{Gr}_K(k,n)\) denote the Grassmannian of \(k\)-dimensional linear subspaces of \(K^n\), where the base field \(K\) remains to be specified. Its Pl\"ucker ideal is
\[
I_{k,n}\subseteq K[p_J: J\subseteq\{1,\ldots,n\},\ |J|=k].
\]
The supplied record refers to the tropicalization of this Pl\"ucker embedding. It does not specify the field, valuation, or fan structure used to interpret its maximal cones.

In a formulation using a trivially valued field and a fan on which initial ideals are constant on relative interiors, a maximal cone \(C\) determines an initial ideal
\[
J_C=\operatorname{in}_{w}(I_{k,n}),\qquad w\in\operatorname{relint}(C).
\]
When the associated initial degeneration is toric, denote it by \(X_C\). A maximal tropical cone does not, merely by being maximal, guarantee a toric initial degeneration. Thus the phrase ``toric degenerations arising from distinct maximal cones'' requires either restricting to cones that give toric degenerations or specifying an additional construction for other cones.

\textbf{Problem, with unresolved qualifications.} For higher Grassmannians, meaning \(k>2\), characterize the combinatorial mutations relating toric degenerations \(X_C\) and \(X_{C'}\) arising from distinct maximal cones \(C\neq C'\) in the tropical Grassmannian. A precise characterization should specify the lattice polytopes or other combinatorial objects attached to these degenerations, the permitted mutation operations, and the conditions under which those objects are related by the permitted operations.

The supplied record does not define ``such mutations'': its quoted original passage refers to preceding context that is not supplied. In particular, it does not determine whether the intended relation is a single combinatorial mutation or a finite sequence of mutations, which polytopes and lattice identifications are intended, or which maximal cones are eligible. The statement therefore remains conditional on these choices. No assertion that every pair of toric degenerations is mutation-related is supplied or made here. The intended range of \(n\), including whether the boundary cases \(n=k\) and \(n=k+1\) are excluded, is also unspecified.
\subsection{Short English statement}
Describe the combinatorial mutations connecting toric degenerations of Grassmannians with \(k>2\) obtained from different maximal tropical cones. The supplied record leaves the mutation convention, associated combinatorial objects, and eligible cones unspecified, so this is a qualified statement requiring mathematical review.
\subsection{Sources}
\begin{itemize}
\item Local registry: \texttt{data/unsolvedmath/problems.json}, permanent UnsolvedMath ID \(30005082\), source problem code \texttt{catalogue entry}, title ``Mutations of Higher Grassmannian Degenerations.'' The supplied record belongs to the Oberwolfach Reports Open Problems collection.
\item Supplied report citation: \emph{Toric Geometry (2023)}, Oberwolfach Reports, pp.\ 861--927, DOI \texttt{10.4171/owr/2022/17}. \url{https://doi.org/10.4171/owr/2022/17}. The record gives proposed year 2022 and names no proposer. These details are retained without resolving the difference between the citation's year and the report identifier.
\item Supplied original passage: ``The natural question is to understand such mutations for higher Grassmannians.'' The antecedent of ``such mutations'' is not included in the supplied record.
\item Supplied literature reference: Clarke--Mohammadi--Zaffalon. \url{https://doi.org/10.1016/j.jalgebra.2023.09.026}. No title, theorem number, or precise supporting passage is supplied.
\item Supplied literature reference: Oberwolfach Report 2022/17. \url{https://doi.org/10.4171/owr/2022/17}.
\item Literature assessment dated exactly \texttt{2026-08-22}: the supplied curation labels the problem open and states that combinatorial mutations are understood for several matching-field families and low-dimensional cases, while the general higher-Grassmannian problem remains open. This is supplied dated curation derived from the OWR material, not an independently verified literature assessment; its status and cited claims require independent verification.
\item The supplied registry labels the statement \texttt{corrected\_verified} and reports reconstruction and checking against the official Oberwolfach report. These are source-provided verification claims, not independent verification performed for this draft.
\end{itemize}
\end{document}
